%% ============================================================
%%  Chapter 3 – Methodology
%% ============================================================
\chapter{Methodology}
\label{ch:method}

\begin{infobox}[Chapter Overview]
This chapter describes the complete \medxplain{} system: the hospital-grade
data infrastructure (DICOM ingestion, PHI scrubbing, report parsing), the
vision-language model architecture, the two-phase training protocol, and the
explainability and clinical validation pipelines.
\end{infobox}

%% ─────────────────────────────────────────────────────────────
\section{System Overview}
\label{sec:system_overview}

Figure~\ref{fig:pipeline} illustrates the end-to-end \medxplain{} pipeline.

\begin{figure}[htbp]
\centering
\begin{tikzpicture}[
  node distance=8mm and 14mm,
  box/.style={rectangle, rounded corners=3pt, draw=blue!60!black,
              fill=blue!8, minimum width=32mm, minimum height=9mm,
              font=\small, text centered},
  redbox/.style={box, draw=red!60!black, fill=red!8},
  greenbox/.style={box, draw=green!50!black, fill=green!6},
  arr/.style={-Stealth, thick, blue!60!black},
]

%% Row 1 – Data layer
\node[box]  (dicom)   {DICOM\\Files (PACS)};
\node[box, right=of dicom]   (loader)  {DICOM\\Loader};
\node[redbox, right=of loader]  (phi)  {PHI\\Scrubber};
\node[greenbox, right=of phi]   (tensor) {Tensor\\Converter};

%% Row 2 – Model layer
\node[box, below=18mm of loader]  (visenc) {BiomedCLIP\\Vision Encoder};
\node[box, right=15mm of visenc]  (fusion) {Cross-Attention\\Fusion};
\node[box, right=15mm of fusion]  (lmdec)  {BioGPT\\+ LoRA};

%% Row 3 – Output layer
\node[greenbox, below=15mm of fusion] (answer)  {Answer +\\Confidence};
\node[redbox,   right=12mm of answer] (gradcam) {Grad-CAM\\Heatmap};
\node[box, left=12mm of answer]       (report)  {Report Parser\\(Optional)};

%% Row 4 – UI
\node[box, below=15mm of answer, xshift=15mm]   (ui) {Gradio Clinical\\Interface};

%% Arrows – data layer
\draw[arr] (dicom)  -- (loader);
\draw[arr] (loader) -- (phi);
\draw[arr] (phi)    -- (tensor);
\draw[arr] (tensor) -- ++(0,-10mm) -| (visenc);

%% Arrows – model layer
\draw[arr] (visenc) -- (fusion);
\draw[arr] (fusion) -- (lmdec);
\draw[arr] (fusion) -- (answer);
\draw[arr] (lmdec)  |- (answer);

%% Arrows – output layer
\draw[arr] (answer)  -- (gradcam);
\draw[arr] (answer)  -- (report);
\draw[arr] (answer)  -- (ui);
\draw[arr] (gradcam) |- (ui);

\end{tikzpicture}
\caption{End-to-end \medxplain{} pipeline. Blue boxes: computation modules;
         red: privacy/explainability; green: data/output.}
\label{fig:pipeline}
\end{figure}

%% ─────────────────────────────────────────────────────────────
\section{Hospital-Grade Data Infrastructure}
\label{sec:data_infra}

\subsection{DICOM Ingestion Pipeline}
\label{sec:dicom}

Hospital imaging data is stored in the Digital Imaging and Communications
in Medicine (\dicom) format, a complex binary standard that encodes pixel
data alongside hundreds of metadata attributes (tags) in a structured
dataset~\citep{dicom2023}.

The \code{DICOMLoader} class (Listing~\ref{lst:dicom}) wraps the
\code{pydicom} library to:
\begin{enumerate}
  \item Recursively discover all \code{.dcm} files in a study directory.
  \item Validate each file against the known transfer syntax whitelist
        (Implicit/Explicit VR, JPEG, JPEG-LS, JPEG~2000 via \code{pylibjpeg}).
  \item Extract safe metadata tags (modality, window centre/width, pixel
        spacing) into the \code{DICOMStudy} dataclass.
  \item Normalise pixel arrays to float32 $\in [0,1]$ via VOI LUT windowing
        (Equation~\ref{eq:window}).
  \item Pseudonymise PatientID and StudyInstanceUID using SHA-256
        (Equation~\ref{eq:pseudo}).
\end{enumerate}

\paragraph{Pixel normalisation.}
The windowing function maps raw Hounsfield units (CT) or grey values
(radiography) to the display range specified by the DICOM
\texttt{WindowCenter} ($C$) and \texttt{WindowWidth} ($W$) tags:
\begin{equation}
  p_{\mathrm{norm}} = \mathrm{clip}\!\left(
    \frac{p - \left(C - W/2\right)}{W},\ 0,\ 1
  \right)
  \label{eq:window}
\end{equation}
When windowing tags are absent, min-max normalisation is applied as a
fallback.

\paragraph{Deterministic pseudonymisation.}
Each identifying value $v$ (PatientID, StudyInstanceUID) is replaced by a
truncated SHA-256 digest of the concatenated application salt and value:
\begin{equation}
  \mathrm{pseudo}(v) = \mathrm{SHA\text{-}256}(\texttt{salt} \| v)[:32]
  \label{eq:pseudo}
\end{equation}
The salt is injected at runtime via the
\code{MEDXPLAIN\_PSEUDO\_SALT} environment variable, ensuring the mapping
is irreversible without the key while preserving longitudinal linkage across
studies from the same patient.

\begin{lstlisting}[caption={Using the DICOM loader.},
                   label={lst:dicom}]
from data_ingestion.dicom_pipeline import DICOMLoader

loader = DICOMLoader(target_size=(224, 224))
study  = loader.load_study("/hospital/pacs/study_001/")

# study.patient_id  ->  "a4f3c2..."  (SHA-256 pseudonym)
# study.modality    ->  "CT"
# study.images[0].shape  ->  torch.Size([1, 224, 224])
\end{lstlisting}

\paragraph{Multi-frame support.}
Single-frame modalities (CR, DX) are returned as tensors of shape
$(1, H, W)$; multi-frame CT/MR volumes are returned as $(N, 1, H, W)$
where $N$ is the slice count.

\subsection{PHI Scrubbing}
\label{sec:phi}

The \code{phi\_scrubber} module provides a one-call API:
\begin{lstlisting}[caption={PHI scrubbing API.}, label={lst:phi}]
from data_ingestion.phi_scrubber import scrub_dicom_study
clean = scrub_dicom_study(raw_study, keep_longitudinal_link=True)
\end{lstlisting}

Internally it performs three passes:

\begin{enumerate}
  \item \textbf{DICOM tag scrub.} All 35+ tags in
        \code{DICOM\_CONFIDENTIALITY\_KEYWORDS} (PatientName, DOB,
        InstitutionName, \dots) are deleted from a deep copy of the dataset.
        Tags in \code{UID\_KEYWORDS\_TO\_PSEUDONYMISE} are replaced with
        deterministic \code{2.25.<hash>} UIDs.
        Tags not in \code{SAFE\_DICOM\_KEYWORDS} are deleted by default.

  \item \textbf{Report text scrub.}
        A battery of compiled regular expressions redacts:
        MRNs (\code{MRN\textbackslash :\textbackslash s*\textbackslash d\{4,12\}}),
        phone numbers, dates (multiple formats), email addresses, and
        physician name signatures (``Reported by Dr.~X'').
        Optionally, a spaCy \textsc{NER} model (\code{en\_core\_web\_sm} or
        \code{en\_core\_sci\_lg}) provides additional PERSON entity
        redaction.

  \item \textbf{Audit record.}
        A single JSON line is appended to
        \code{./logs/phi\_audit.log} containing: ISO-8601 timestamps,
        SHA-256 hash of the source file paths (never the content), frame
        count, and a boolean flag indicating whether NER was used.
        This provides a tamper-evident compliance record without logging PHI.
\end{enumerate}

\subsection{Radiology Report Parser}
\label{sec:report_parser}

Radiology reports follow a semi-structured free-text format with named
sections (Clinical History, Findings, Impression). The
\code{ReportParser} class extracts these sections using line-by-line header
matching and then applies the 14-entry \textbf{pathology catalogue}
(Table~\ref{tab:pathology}) to the combined Findings + Impression text.

\begin{table}[htbp]
\centering
\caption{Pathology catalogue used by \code{ReportParser}.}
\label{tab:pathology}
\begin{tabular}{@{}llp{6cm}@{}}
\toprule
\textbf{Key} & \textbf{Display Name} & \textbf{Detection Patterns} \\
\midrule
\code{pneumothorax}      & Pneumothorax          & pneumothorax, ptx \\
\code{pleural\_effusion} & Pleural Effusion       & pleural effusion, hydrothorax \\
\code{pneumonia}         & Pneumonia              & pneumonia, pneumonitis, consolidation \\
\code{atelectasis}       & Atelectasis            & atelectasis, collapse \\
\code{cardiomegaly}      & Cardiomegaly           & cardiomegaly, cardiac enlargement \\
\code{pulmonary\_edema}  & Pulmonary Oedema       & pulmonary oedema, vascular congestion \\
\code{nodule}            & Pulmonary Nodule       & nodule, rounded opacity \\
\code{fracture}          & Fracture               & fracture, cortical disruption \\
\code{pneumoperitoneum}  & Pneumoperitoneum       & pneumoperitoneum, free air \\
\code{pleural\_thickening}& Pleural Thickening    & pleural thickening \\
\code{pleural\_mass}     & Pleural Mass           & pleural mass, pleural lesion \\
\code{hilar\_enlargement}& Hilar Enlargement      & hilar enlargement \\
\code{interstitial}      & Interstitial Pattern   & interstitial markings, reticular \\
\code{aortic\_widening}  & Aortic Widening        & aortic widening, aortic aneurysm \\
\bottomrule
\end{tabular}
\end{table}

For each finding, negation heuristics check the eight words preceding the
detected keyword for terms such as \emph{no}, \emph{without},
\emph{absence of}, or \emph{ruled out}. Laterality (left/right/bilateral)
and severity (mild/moderate/severe) qualifiers are extracted from a 60-word
window surrounding each match.

%% ─────────────────────────────────────────────────────────────
\section{Model Architecture}
\label{sec:architecture}

Figure~\ref{fig:arch} details the three-component architecture.

\begin{figure}[htbp]
\centering
\begin{tikzpicture}[
  node distance=6mm and 18mm,
  block/.style={rectangle, rounded corners=4pt, draw=blue!60!black,
               fill=blue!7, minimum width=36mm, minimum height=14mm,
               text centered, font=\small},
  lora/.style={block, draw=orange!70, fill=orange!8},
  arr/.style={-Stealth, thick},
  label/.style={font=\tiny, text=gray},
]

\node[block] (img) {Input Image\\$224 \times 224$};
\node[block, right=22mm of img] (vit)
      {BiomedCLIP\\ViT-B/16\\{\tiny $\sim$86M params}};
\node[block, right=22mm of vit] (xattn)
      {Cross-Attention\\Fusion\\{\tiny $\sim$9M params, trainable}};
\node[lora,  right=22mm of xattn] (biogpt)
      {BioGPT\\+ LoRA (r=8)\\{\tiny $\sim$2M trainable}};
\node[block, below=14mm of xattn] (head)
      {Answer Head\\$\mathbf{W} \in \mathbb{R}^{768 \times |\mathcal{V}|}$};
\node[block, left=18mm of head] (question) {Question\\Tokens};

\draw[arr, blue!60] (img)      -- (vit);
\draw[arr, blue!60] (vit)      -- node[above, label]{visual tokens} (xattn);
\draw[arr, blue!60] (question) -- (xattn);
\draw[arr, blue!60] (xattn)   -- (biogpt);
\draw[arr, blue!60] (biogpt)  -| (head);
\draw[arr, green!50!black] (head) -- ++(0,-12mm)
  node[below, font=\small]{Answer logits};

%% Freeze / unfreeze annotations
\node[label, above=1mm of vit] {Phase 1: frozen\\Phase 2: last 2 blocks};
\node[label, above=1mm of biogpt] {Phase 1: frozen\\Phase 2: LoRA unfrozen};
\end{tikzpicture}
\caption{\medxplain{} model architecture. Orange = LoRA-adapted layers.
         Arrows show data flow. Freeze/unfreeze status is shown per training phase.}
\label{fig:arch}
\end{figure}

\paragraph{Vision encoder.}
BiomedCLIP~\citep{zhang2023biomedclip} encodes input images of size
$224 \times 224$ into a sequence of patch tokens using its ViT-B/16
backbone. The \texttt{[CLS]} token and patch tokens are passed to the
fusion layer as visual tokens $\mathbf{V} \in \mathbb{R}^{N_v \times D}$
($N_v = 197$, $D = 768$).

\paragraph{Cross-attention fusion.}
The fusion module implements standard multi-head attention where text
features act as queries and visual features as keys and values:
\begin{equation}
  \mathrm{Fusion}(\mathbf{Q}, \mathbf{K}, \mathbf{V}) =
  \mathrm{softmax}\!\left(\frac{\mathbf{Q}\mathbf{K}^\top}{\sqrt{D}}\right)
  \mathbf{V}
  \label{eq:crossattn}
\end{equation}
with 8~heads and $D=768$. Implemented in
\code{medical\_vqa\_infrastructure/models/fusion.py}.

\paragraph{Language decoder.}
BioGPT processes the fused representation conditioned on the question
tokens. LoRA adapters (Equation~\ref{eq:lora}) are injected into the
query and value projections of every attention layer.

\paragraph{Answer head.}
A linear projection maps the pooled fused feature vector to a
classification over the closed answer vocabulary
$\mathcal{V}$ (e.g., \{Yes, No\}).

%% ─────────────────────────────────────────────────────────────
\section{Two-Phase Training}
\label{sec:training}

Training is implemented in
\code{training/train\_vqa\_two\_phase.py} and is launched as:
\begin{lstlisting}[language=bash, caption={Launching two-phase training.},
                   label={lst:train}]
python training/train_vqa_two_phase.py \
    --config config/hospital_config.yaml \
    --phase1_epochs 10 --phase2_epochs 20 \
    --batch_size 8 --amp
\end{lstlisting}

\subsection{Phase 1: Vision-Language Alignment}
\label{subsec:phase1}

\paragraph{Objective.}
Learn a shared embedding space such that images and their corresponding
clinical questions are proximate.

\paragraph{Trainable parameters.}
The vision encoder and language decoder are \emph{frozen}. Only the
cross-attention fusion layer and two linear projection heads
($\mathbf{W}_{\mathrm{img}}, \mathbf{W}_{\mathrm{txt}} \in
\mathbb{R}^{D \times D}$) receive gradient updates.

\paragraph{Loss function.}
The NT-Xent (InfoNCE) contrastive loss (Equation~\ref{eq:ntxent}) is applied
symmetrically over the batch:
\begin{equation}
  \mathcal{L}_1 = \frac{1}{2}\bigl(
    \mathcal{L}_{\mathrm{i2t}} + \mathcal{L}_{\mathrm{t2i}}
  \bigr)
  \label{eq:phase1_loss}
\end{equation}
with temperature $\tau = 0.07$.

\paragraph{Optimiser.}
AdamW~\citep{loshchilov2019adamw} with learning rate $\eta_1 = 10^{-4}$
and weight decay $10^{-4}$, combined with cosine annealing over 10 epochs.

\paragraph{Checkpoint.}
\code{checkpoints/phase1\_alignment.pt}

\subsection{Phase 2: VQA Fine-tuning}
\label{subsec:phase2}

\paragraph{Objective.}
Answer closed-form clinical questions accurately.

\paragraph{Initialisation.}
Phase~1 checkpoint is loaded. The entire model is frozen, then:
\begin{itemize}
  \item The last 2 transformer blocks of the vision encoder are unfrozen.
  \item LoRA adapters (rank~$r=8$, $\alpha=16$) are injected into the
        \texttt{query} and \texttt{value} projections of BioGPT.
  \item The cross-attention fusion and answer head are unfrozen.
\end{itemize}

This strategy protects the low-level visual representations learned during
Phase~1 while allowing the model to specialise the higher-level features
for the VQA task.

\paragraph{Loss function.}
Standard cross-entropy over the answer vocabulary $\mathcal{V}$:
\begin{equation}
  \mathcal{L}_2 = -\sum_{c \in \mathcal{V}} y_c \log \hat{p}_c
  \label{eq:ce}
\end{equation}

\paragraph{Optimiser.}
AdamW with $\eta_2 = 5 \times 10^{-5}$, gradient clipping
$\|\nabla\|_2 \leq 1.0$, cosine annealing over 20 epochs.

\paragraph{Regularisation.}
Automatic Mixed Precision (AMP, \code{torch.cuda.amp}) halves GPU memory
consumption. Gradient checkpointing is optionally enabled for
memory-constrained hardware.

\paragraph{Checkpoint.}
\code{checkpoints/phase2\_vqa\_finetuned.pt}

\subsection{Training Infrastructure}
\label{subsec:train_infra}

\begin{itemize}
  \item \textbf{WandB logging}: Loss curves, learning rates, and validation
        accuracy are logged every step. Graceful fallback to stdout if WandB
        is unavailable.
  \item \textbf{Validation}: Every 50 gradient steps, the model is evaluated
        on the held-out validation set. Sample predictions are logged for
        qualitative monitoring.
  \item \textbf{Windows compatibility}: \code{num\_workers=0} (main-process
        DataLoader) avoids the multiprocessing \texttt{fork} issue on Windows.
\end{itemize}

%% ─────────────────────────────────────────────────────────────
\section{Explainability Integration}
\label{sec:explainability}

\subsection{Grad-CAM Computation}
\label{subsec:gradcam_impl}

For the VQA setting, the answer class $c$ corresponds to the argmax of the
answer head logits. Grad-CAM is computed with respect to the feature maps
of the \emph{last} attention block of the BiomedCLIP ViT, following
the formulation in Equations~\ref{eq:alpha}--\ref{eq:gradcam}.

The resulting heatmap is:
\begin{enumerate}
  \item Resized to the input resolution ($224 \times 224$) using bilinear
        interpolation.
  \item Converted to a jet colourmap ($\in [0,1]$).
  \item Alpha-blended with the original image at $\alpha = 0.5$.
\end{enumerate}
Implementation: \code{explainability/grad\_cam.py} (via
\code{pytorch\_grad\_cam}).

\subsection{Additional Attribution Methods}
\label{subsec:xai_additional}

\begin{itemize}
  \item \textbf{Integrated Gradients}~(\code{explainability/integrated\_gradients.py}):
        Provides axiomatic attribution with completeness guarantee.
        Used for academic analysis of model behaviour.
  \item \textbf{Counterfactual Explanations}~(\code{explainability/counterfactual.py}):
        Identifies the minimal pixel-space perturbation that flips the
        predicted answer. Intended for offline discrepancy case analysis.
\end{itemize}

%% ─────────────────────────────────────────────────────────────
\section{Clinical User Interface}
\label{sec:ui}

The \code{vqa\_app\_deliverable/app\_gradio.py} application implements
the Cahier des Charges UI specification using Gradio Blocks:

\begin{itemize}
  \item \textbf{Privacy banner}: A persistent markdown block at the top of
        every tab reads: ``\texttt{LOCAL PROCESSING ONLY -- No data leaves
        this machine}''.
  \item \textbf{DICOM upload}: \code{gr.File(file\_types=[".dcm",".png",".jpg"])}
        accepts DICOM and standard raster formats. On upload, the file is
        decoded by \code{DICOMLoader} and converted to a PIL image for display.
  \item \textbf{Two-panel output}: The VQA answer and confidence score are
        shown on the left; the Grad-CAM heatmap is rendered on the right
        at the same resolution as the uploaded image.
  \item \textbf{DICOM metadata panel}: A \code{gr.Accordion} containing a
        \code{gr.JSON} widget displays the pseudonymised study metadata.
  \item \textbf{Clinical tabs}: Report-Aware, Context-Aware, Longitudinal
        View, Differential Diagnosis, and One-Click Report tabs are
        accessible alongside the primary VQA tab.
\end{itemize}

%% ─────────────────────────────────────────────────────────────
\section{Clinical Validation}
\label{sec:clin_val_method}

The \code{evaluation/clinical\_validation.py} module implements
radiologist-comparison validation:

\begin{enumerate}
  \item Load a CSV of ground-truth annotations
        (\texttt{image\_id}, \texttt{question}, \texttt{radiologist\_answer}).
  \item For each row, run model inference and compare to the radiologist label.
  \item Compute Cohen's Kappa $\kappa$ and exact-match accuracy.
  \item For every discrepancy, generate a PHI-free PDF report containing
        the source image, question, model answer and confidence,
        radiologist answer, and Grad-CAM overlay.
\end{enumerate}

Cohen's Kappa is defined as:
\begin{equation}
  \kappa = \frac{p_o - p_e}{1 - p_e}
  \label{eq:kappa}
\end{equation}
where $p_o$ is the observed agreement and $p_e$ is the expected agreement
under independence.
